\section{Introduction}
\label{sec:introduction}

If a model recognises that it is being evaluated, its behaviour under evaluation may not
predict its behaviour in deployment. Researchers have built detectors to flag this evaluation
awareness. Some read only text. For instance, one detector shows a model a transcript
and asks whether it comes from an evaluation \citep{needham2025large}. Others read the model's
internals, as an activation probe does \citep{nguyen2025probing}. If these detectors measure the
same property, they should flag the same responses.

One check asks whether a detector separates groups of prompts known to differ. Psychometrics
calls this known-groups validity: groups expected to differ in a property should also differ in
its measure \citep{cronbach1955construct,hattie1984procedures}. EvalAwareBench
\citep{li2026decomposing} supplies such groups. It adds test-like cues such as placeholder names
to otherwise unchanged prompts. A detector passes the check if it flags these cued prompts more
often than the plain ones.

However, this check is weak. Every detector we analysed passes it, but so do random directions in
Qwen's activations after the same tuning as the activation probe. These directions
were drawn without reference to awareness, so a check that they pass cannot tell good detectors
from bad. \citet{grover2026can} report a similar result for steering: some random directions
shift a model's behaviour more than the learned direction does (other related work in
Appendix~\ref{app:related}).

A sharper check tests convergent validity, which asks whether different methods agree on the
same property \citep{campbell1959convergent}. Over all prompts, though, the cues themselves
create some agreement. Because every detector that responds to the cues flags cued prompts more
often, two detectors can agree overall even if they flag independently among prompts with the
same cues. We therefore compare detectors only among such prompts.

We analysed five detectors on Qwen 3.6 27B and three on Gemma 4 31B
\citep{qwenteam2026qwen36,gemmateam2026gemma4}. Our contributions are:
\begin{itemize}
  \item We show that passing a known-groups check might not be enough. On Qwen, random activation
  directions pass it, and so does a contrastive activation probe that, among prompts with the same cues, agrees
  with none of the three other detectors it can be compared with.
  \item We compare detectors only among prompts with the same cues (post hoc) and against the
  highest Cohen's $\kappa$ their positive rates allow. This separates flagging at different rates from
  flagging different responses.
  \item We find that, among prompts with the same cues (post hoc), self-report and
  third-person reach most of the agreement their positive rates allow, and that most of their
  disagreement comes from flagging at different rates. Because both are the same model judging
  the same response, this is an upper bound on agreement across methods. In contrast, the Qwen activation probe
  agrees with them no more than chance.
\end{itemize}
